
\documentclass[11pt,a4paper]{article}
\usepackage[T1]{fontenc}
\usepackage{lmodern}
\usepackage{amsmath,amssymb,amsfonts,mathtools}
\usepackage{graphicx,booktabs,longtable,array,tabularx,adjustbox}
\usepackage{geometry,microtype,fancyhdr,lastpage,etoolbox,enumitem}
\usepackage{newunicodechar,textcomp}
\usepackage{hyperref}
\hypersetup{hidelinks,pdfcreator={XeLaTeX / QM1 academic edition}}
\geometry{left=20mm,right=20mm,top=22mm,bottom=23mm}
\setcounter{secnumdepth}{2}
\setlength{\parindent}{0pt}
\setlength{\parskip}{4pt plus 1pt}
\setlength{\emergencystretch}{2em}
\setlength{\tabcolsep}{3pt}
\renewcommand{\arraystretch}{1.16}
\setlist{nosep,topsep=3pt,itemsep=2pt,leftmargin=7mm}
\AtBeginEnvironment{longtable}{\footnotesize\setlength{\tabcolsep}{2.2pt}}
\newcommand{\tightlist}{\setlength{\itemsep}{0pt}\setlength{\parskip}{0pt}}
\providecommand{\Cov}{\operatorname{Cov}}
\providecommand{\Var}{\operatorname{Var}}
\providecommand{\Corr}{\operatorname{Corr}}
\providecommand{\MAD}{\operatorname{MAD}}
\providecommand{\E}{\mathbb{E}}
\providecommand{\R}{\mathbb{R}}
\newunicodechar{—}{\textemdash}
\newunicodechar{–}{\textendash}
\newunicodechar{€}{\texteuro}
\newunicodechar{≤}{\ensuremath{\leq}}
\newunicodechar{≥}{\ensuremath{\geq}}
\newunicodechar{×}{\ensuremath{\times}}
\newunicodechar{÷}{\ensuremath{\div}}
\newunicodechar{Σ}{\ensuremath{\Sigma}}
\newunicodechar{σ}{\ensuremath{\sigma}}
\newunicodechar{μ}{\ensuremath{\mu}}
\newunicodechar{Δ}{\ensuremath{\Delta}}
\newunicodechar{∑}{\ensuremath{\sum}}
\newunicodechar{≈}{\ensuremath{\approx}}
\newunicodechar{√}{\ensuremath{\surd}}
\newunicodechar{≠}{\ensuremath{\ne}}
\newunicodechar{→}{\ensuremath{\to}}
\newunicodechar{∞}{\ensuremath{\infty}}
\newunicodechar{∈}{\ensuremath{\in}}
\newunicodechar{−}{\ensuremath{-}}
\newunicodechar{²}{\ensuremath{^{2}}}
\newunicodechar{³}{\ensuremath{^{3}}}
\newunicodechar{₀}{\ensuremath{_{0}}}
\newunicodechar{₁}{\ensuremath{_{1}}}
\newunicodechar{₂}{\ensuremath{_{2}}}
\newunicodechar{₃}{\ensuremath{_{3}}}
\newunicodechar{₄}{\ensuremath{_{4}}}
\newunicodechar{ₜ}{\ensuremath{_{t}}}
\newunicodechar{ₚ}{\ensuremath{_{p}}}
\newunicodechar{ₙ}{\ensuremath{_{n}}}
\newunicodechar{ₓ}{\ensuremath{_{x}}}
\newunicodechar{ᵢ}{\ensuremath{_{i}}}
\newunicodechar{ŷ}{\ensuremath{\hat y}}
\newunicodechar{Ŷ}{\ensuremath{\widehat Y}}
\newunicodechar{ₘ}{\ensuremath{_{m}}}
\newunicodechar{ₑ}{\ensuremath{_{e}}}
\newunicodechar{ₐ}{\ensuremath{_{a}}}
\newunicodechar{ₖ}{\ensuremath{_{k}}}
\newunicodechar{ₛ}{\ensuremath{_{s}}}
\newunicodechar{ₗ}{\ensuremath{_{l}}}
\newunicodechar{ₕ}{\ensuremath{_{h}}}
\newunicodechar{ₒ}{\ensuremath{_{o}}}
\newunicodechar{₊}{\ensuremath{_{+}}}
\newunicodechar{₋}{\ensuremath{_{-}}}
\newunicodechar{₌}{\ensuremath{_{=}}}
\newunicodechar{₍}{\ensuremath{_{(}}}
\newunicodechar{₎}{\ensuremath{_{)}}}
\newunicodechar{⁻}{\ensuremath{^{-}}}
\newunicodechar{⁺}{\ensuremath{^{+}}}
\newunicodechar{⁰}{\ensuremath{^{0}}}
\newunicodechar{¹}{\ensuremath{^{1}}}
\newunicodechar{⁴}{\ensuremath{^{4}}}
\newunicodechar{⁵}{\ensuremath{^{5}}}
\newunicodechar{⁶}{\ensuremath{^{6}}}
\newunicodechar{⁷}{\ensuremath{^{7}}}
\newunicodechar{⁸}{\ensuremath{^{8}}}
\newunicodechar{⁹}{\ensuremath{^{9}}}
\makeatletter
\renewcommand{\section}{\@startsection{section}{1}{0pt}{2.0ex plus .7ex minus .3ex}{1.0ex}{\large\bfseries}}
\renewcommand{\subsection}{\@startsection{subsection}{2}{0pt}{1.8ex plus .5ex minus .2ex}{.7ex}{\normalsize\bfseries}}
\renewcommand{\subsubsection}{\@startsection{subsubsection}{3}{0pt}{1.3ex plus .4ex minus .2ex}{.5ex}{\normalsize\bfseries}}
\makeatother
\pagestyle{fancy}
\fancyhf{}
\renewcommand{\headrulewidth}{0.35pt}
\renewcommand{\footrulewidth}{0pt}
\fancyhead[L]{\small\textsc{KEDGE Business School --- QM1}}
\fancyhead[R]{\small\textsc{2026--2027}}
\fancyfoot[L]{\small Badr ABOUFARASSE, PhD}
\fancyfoot[R]{\small\thepage\ /\ \pageref{LastPage}}
\setlength{\headheight}{14pt}
\setkeys{Gin}{keepaspectratio}
\begin{document}

\begin{center}
{\small\textsc{KEDGE Business School \quad EBP1 \quad Quantitative Methods 1}}\\[6pt]
{\Large\bfseries Fiche 6 — Nombres-indices}\\[6pt]
{\normalsize\bfseries Corrigés intégralement détaillés}\\[6pt]
{\small Badr ABOUFARASSE, PhD \quad --- \quad 2026--2027}
\end{center}
\vspace{2mm}\hrule\vspace{6mm}

\hypertarget{applications-fondamentales}{%
\section{Applications fondamentales}\label{applications-fondamentales}}

Maîtriser les définitions et les calculs essentiels du cours.

\emph{Application du cours --- 20 min}

\hypertarget{exercice-1-indices-uxe9luxe9mentaires-rebasage-et-propriuxe9tuxe9s}{%
\subsection{Exercice 1 -- Indices élémentaires, rebasage et propriétés}\label{exercice-1-indices-uxe9luxe9mentaires-rebasage-et-propriuxe9tuxe9s}}

Un service coûte 80 euros en 2023, 92 euros en 2024 et 101.20 euros en 2025.

\begin{enumerate}
\item
  Calculer les indices 2024/2023, 2025/2024 et 2025/2023, base 100.
\item
  Vérifier la circularité.
\item
  Calculer l'indice inverse 2023/2025 et vérifier la réversibilité.
\item
  Rebaser la série pour prendre 2024 comme base 100.
\item
  Interpréter correctement l'indice 126.5.
\end{enumerate}

\hypertarget{corriguxe9-duxe9tailluxe9}{%
\subsubsection{Corrigé détaillé}\label{corriguxe9-duxe9tailluxe9}}

{\[I_{2024/2023}=100\frac{92}{80}=115.\]} {\[I_{2025/2024}=100\frac{101.20}{92}=110.\]} {\[I_{2025/2023}=100\frac{101.20}{80}=126.5.\]} Circularité : {\[\frac{115\times110}{100}=126.5.\]} Indice inverse : {\[I_{2023/2025}=100\frac{80}{101.20}\approx79.0514.\]} Réversibilité : {\[126.5\times79.0514\approx10000.\]} Rebasage en 2024=100 : {\[I_{2023}^{(2024)}=100\frac{100}{115}\approx86.96,\]} {\[I_{2024}^{(2024)}=100,
\qquad
I_{2025}^{(2024)}=110.\]} L'indice 126.5 signifie que le niveau de 2025 est 26.5\% supérieur à celui de 2023 ; il ne signifie pas une hausse de 126.5\%.

\emph{Application du cours --- 25 min}

\hypertarget{compluxe9ment-1-duxe9composition-compluxe8te-prixquantituxe9s}{%
\subsection{Complément 1 --- Décomposition complète prix--quantités}\label{compluxe9ment-1-duxe9composition-compluxe8te-prixquantituxe9s}}

Un service achats observe :

\begin{longtable}[]{@{}lllll@{}}
\toprule\noalign{}
Produit & p0 & q0 & p1 & q1 \\
\midrule\noalign{}
\endhead
\bottomrule\noalign{}
\endlastfoot
A & 4 & 100 & 5 & 90 \\
B & 10 & 50 & 11 & 60 \\
C & 20 & 20 & 18 & 30 \\
\end{longtable}

\begin{enumerate}
\tightlist
\item
  Calculez les quatre totaux Σp0q0, Σp1q0, Σp0q1 et Σp1q1.
\item
  Calculez les indices de valeur, de prix de Laspeyres, Paasche et Fisher.
\item
  Interprétez l'écart entre valeur et prix.
\end{enumerate}

\textbf{Indication méthodologique :}

Complétez les quatre colonnes de valeurs croisées avant tout rapport.

\hypertarget{corriguxe9-duxe9tailluxe9-1}{%
\subsubsection{Corrigé détaillé}\label{corriguxe9-duxe9tailluxe9-1}}

\textbf{Totaux.} Σp0q0 = \textbf{1 300} ; Σp1q0 = \textbf{1 410} ; Σp0q1 = \textbf{1 560} ; Σp1q1 = \textbf{1 650}.

\textbf{Indice de valeur.} 1 650/1 300 × 100 = \textbf{126,92}.

\textbf{Laspeyres-prix.} \textbf{108,46}. \textbf{Paasche-prix.} \textbf{105,77}. \textbf{Fisher-prix.} √(108,46 × 105,77) = \textbf{107,11}.

\textbf{Interprétation.} La valeur des achats augmente de 26,92 \%, alors que la hausse équilibrée des prix est d'environ 7,11 \%. L'essentiel de la hausse de valeur provient donc des quantités ou de la composition du panier, sans que les indices expliquent la cause de ce changement.

Détail de chacun des quatre agrégats (à vérifier produit par produit) :

\begin{longtable}[]{@{}lllll@{}}
\toprule\noalign{}
Produit / Product & p₀q₀ & p₁q₀ & p₀q₁ & p₁q₁ \\
\midrule\noalign{}
\endhead
\bottomrule\noalign{}
\endlastfoot
A & 400 & 500 & 360 & 450 \\
B & 500 & 550 & 600 & 660 \\
C & 400 & 360 & 600 & 540 \\
Σ & 1300 & 1410 & 1560 & 1650 \\
\end{longtable}

Laspeyres = 1410/1300 × 100 = 108.4615 ; Paasche = 1650/1560 × 100 = 105.7692 ; Fisher = √(108.4615 × 105.7692) = 107.1069.

\hypertarget{pruxe9paration-aux-uxe9valuations}{%
\section{Préparation aux évaluations}\label{pruxe9paration-aux-uxe9valuations}}

Résoudre des exercices de contrôle continu et de partiel en autonomie.

\emph{Préparation CC / partiel --- 35 min}

\hypertarget{exercice-2-laspeyres-paasche-fisher-et-factor-reversal}{%
\subsection{Exercice 2 -- Laspeyres, Paasche, Fisher et factor reversal}\label{exercice-2-laspeyres-paasche-fisher-et-factor-reversal}}

\begin{longtable}[]{@{}lllll@{}}
\toprule\noalign{}
& {\(p_0\)} & {\(q_0\)} & {\(p_1\)} & {\(q_1\)} \\
\midrule\noalign{}
\endhead
\bottomrule\noalign{}
\endlastfoot
Produit X & 4 & 100 & 5 & 90 \\
Produit Y & 10 & 40 & 12 & 50 \\
\end{longtable}

\begin{enumerate}
\item
  Calculer les quatre agrégats {\(\sum p_0q_0\)}, {\(\sum p_1q_0\)}, {\(\sum p_0q_1\)}, {\(\sum p_1q_1\)}.
\item
  Calculer les indices de prix de Laspeyres, Paasche et Fisher.
\item
  Calculer les indices de quantité de Laspeyres, Paasche et Fisher.
\item
  Calculer l'indice de valeur.
\item
  Vérifier la propriété de factor reversal de Fisher sur une échelle base 100.
\end{enumerate}

\hypertarget{corriguxe9-duxe9tailluxe9-2}{%
\subsubsection{Corrigé détaillé}\label{corriguxe9-duxe9tailluxe9-2}}

Les agrégats sont : {\[\sum
p_0q_0=4(100)+10(40)=800,\]} {\[\sum
p_1q_0=5(100)+12(40)=980,\]} {\[\sum
p_0q_1=4(90)+10(50)=860,\]} {\[\sum
p_1q_1=5(90)+12(50)=1050.\]} Prix : {\[I_P^L=100\frac{980}{800}=122.50,\]} {\[I_P^P=100\frac{1050}{860}\approx122.09,\]} {\[I_P^F=\sqrt{122.50\times122.09}\approx122.30.\]} Quantités : {\[I_Q^L=100\frac{860}{800}=107.50,\]} {\[I_Q^P=100\frac{1050}{980}\approx107.14,\]} {\[I_Q^F=\sqrt{107.50\times107.14}\approx107.32.\]} Valeur : {\[I_V=100\frac{1050}{800}=131.25.\]} Factor reversal : {\[\frac{I_P^F I_Q^F}{100}
\approx\frac{122.2963\times107.3213}{100}
=131.25=I_V.\]}

\hypertarget{vuxe9rification-des-quatre-agruxe9gats-produit-par-produit}{%
\subsubsection{Vérification des quatre agrégats produit par produit}\label{vuxe9rification-des-quatre-agruxe9gats-produit-par-produit}}

\begin{longtable}[]{@{}lllll@{}}
\toprule\noalign{}
Produit & p₀q₀ & p₁q₀ & p₀q₁ & p₁q₁ \\
\midrule\noalign{}
\endhead
\bottomrule\noalign{}
\endlastfoot
X & 4×100=400 & 5×100=500 & 4×90=360 & 5×90=450 \\
Y & 10×40=400 & 12×40=480 & 10×50=500 & 12×50=600 \\
Σ & 800 & 980 & 860 & 1050 \\
\end{longtable}

\[\begin{gathered}
L_P=100\frac{980}{800},\\P_P=100\frac{1050}{860},\\L_Q=100\frac{860}{800},\\P_Q=100\frac{1050}{980}.
\end{gathered}\]

Vérifier la propriété de Fisher sur les rapports non arrondis : Fp×Fq/100=100×1050/800=131,25.

\emph{Préparation CC / partiel --- 45 min}

\hypertarget{exercice-4-indices-complets-sur-quatre-produits-type-partiel}{%
\subsection{Exercice 4 --- Indices complets sur quatre produits --- type partiel}\label{exercice-4-indices-complets-sur-quatre-produits-type-partiel}}

Pour quatre produits, les prix et quantités sont observés à deux dates. Base 0 = 100.

\begin{longtable}[]{@{}lllll@{}}
\toprule\noalign{}
Produit / Product & p₀ & q₀ & p₁ & q₁ \\
\midrule\noalign{}
\endhead
\bottomrule\noalign{}
\endlastfoot
A & 53 & 20 & 62 & 25 \\
B & 41 & 28 & 52 & 34 \\
C & 75 & 35 & 82 & 38 \\
D & 25 & 30 & 38 & 33 \\
\end{longtable}

\begin{enumerate}
\tightlist
\item
  Construire les quatre colonnes de produits p₀q₀, p₁q₀, p₀q₁, p₁q₁ et leurs totaux.
\item
  Calculer les indices de prix et de quantité de Laspeyres, Paasche et Fisher.
\item
  Calculer l\textquotesingle indice de valeur et vérifier la propriété de Fisher. Interpréter les trois indices.
\end{enumerate}

\hypertarget{corriguxe9-duxe9tailluxe9-3}{%
\subsubsection{Corrigé détaillé}\label{corriguxe9-duxe9tailluxe9-3}}

\begin{longtable}[]{@{}lllll@{}}
\toprule\noalign{}
Produit / Product & p₀q₀ & p₁q₀ & p₀q₁ & p₁q₁ \\
\midrule\noalign{}
\endhead
\bottomrule\noalign{}
\endlastfoot
A & 1060 & 1240 & 1325 & 1550 \\
B & 1148 & 1456 & 1394 & 1768 \\
C & 2625 & 2870 & 2850 & 3116 \\
D & 750 & 1140 & 825 & 1254 \\
Σ & 5583 & 6706 & 6394 & 7688 \\
\end{longtable}

Les quatre agrégats sont A=Σp₀q₀=5583, B=Σp₁q₀=6706, C=Σp₀q₁=6394, D=Σp₁q₁=7688.

\[L_P=100B/A=100(6706/5583)=120.114634,\quad P_P=100D/C=100(7688/6394)=120.237723.\]

\[L_Q=100C/A=100(6394/5583)=114.526240,\quad P_Q=100D/B=100(7688/6706)=114.643603.\]

\[F_P=\sqrt{L_P P_P}=120.176163,\quad F_Q=\sqrt{L_Q P_Q}=114.584907.\]

\[I_V=100D/A=100(7688/5583)=137.703744,\quad\frac{F_PF_Q}{100}=137.703744.\]

Les prix augmentent d\textquotesingle environ 20,18 \%, les quantités d\textquotesingle environ 14,58 \%, et la valeur totale de 37,70 \%. Ne pas additionner les variations de prix et de quantités : leur composition est multiplicative.

\hypertarget{approfondissement-facultatif}{%
\section{Approfondissement facultatif}\label{approfondissement-facultatif}}

Aller plus loin dans l'interprétation et les applications professionnelles.

\emph{Approfondissement --- 15 min}

\hypertarget{exercice-3-indices-chauxeenuxe9s-et-duxe9flation}{%
\subsection{Exercice 3 -- Indices chaînés et déflation}\label{exercice-3-indices-chauxeenuxe9s-et-duxe9flation}}

Les indices annuels chaînés sont 105, 108 et 97.

\begin{enumerate}
\item
  Calculer l'indice cumulé sur les trois périodes.
\item
  Interpréter le résultat et expliquer pourquoi il ne faut pas additionner les variations annuelles.
\item
  Un chiffre d'affaires nominal vaut 2.4 millions d'euros lorsque l'indice des prix vaut 120. Calculer la valeur réelle en prix de la période de base.
\end{enumerate}

\hypertarget{corriguxe9-duxe9tailluxe9-4}{%
\subsubsection{Corrigé détaillé}\label{corriguxe9-duxe9tailluxe9-4}}

{\[I_{3/0}^{chain}=100\times1.05\times1.08\times0.97.\]} {\[I_{3/0}^{chain}=109.998\approx110.\]} Le niveau final est donc environ 10\% supérieur au niveau initial, malgré la baisse de 3\% de la dernière période. Les taux successifs se composent multiplicativement et ne s'additionnent pas.

Déflation : {\[\text{Valeur réelle}=\frac{2.4}{120/100}=\frac{2.4}{1.20}=2.0\text{ millions
d'euros}.\]}

\hypertarget{tableau-des-facteurs-chauxeenuxe9s}{%
\subsubsection{Tableau des facteurs chaînés}\label{tableau-des-facteurs-chauxeenuxe9s}}

\begin{longtable}[]{@{}lll@{}}
\toprule\noalign{}
Période & Facteur & Indice cumulé \\
\midrule\noalign{}
\endhead
\bottomrule\noalign{}
\endlastfoot
0 & 1 & 100 \\
1 & 1.05 & 105 \\
2 & 1.05×1.08=1.134 & 113.4 \\
3 & 1.134×0.97=1.09998 & 109.998 \\
\end{longtable}

\[I_{3/0}=100(1.05)(1.08)(0.97)=109.998.\]

La somme des pourcentages (+5+8−3=10) ne donne qu'une approximation fortuite du taux composé exact (+9,998 \%). Le déflateur 120 signifie un facteur de prix 1,20 : 2,4/1,20=2,0 millions d'euros réels.

\emph{Approfondissement --- 25 min}

\hypertarget{compluxe9ment-2-changement-de-base-inflation-et-performance-ruxe9elle}{%
\subsection{Complément 2 --- Changement de base, inflation et performance réelle}\label{compluxe9ment-2-changement-de-base-inflation-et-performance-ruxe9elle}}

Le chiffre d'affaires nominal vaut 500 k€, 560 k€ et 630 k€. L'indice des prix correspondant (année 1 = 100) vaut 100, 108 et 120.

\begin{enumerate}
\tightlist
\item
  Calculez le chiffre d'affaires réel en euros de l'année 1.
\item
  Calculez les croissances nominale et réelle entre les années 1 et 3.
\item
  Rebasez l'indice des prix avec année 2 = 100.
\item
  Répondez au responsable qui affirme que l'activité a progressé de 26 \%.
\end{enumerate}

\textbf{Indication méthodologique :}

Déflatez en divisant la valeur nominale par indice/100. Le changement de base ne modifie pas la trajectoire d'inflation.

\hypertarget{corriguxe9-duxe9tailluxe9-5}{%
\subsubsection{Corrigé détaillé}\label{corriguxe9-duxe9tailluxe9-5}}

\textbf{Chiffre d'affaires réel.} Année 1 : \textbf{500 k€} ; année 2 : 560/1,08 = \textbf{518,52 k€} ; année 3 : 630/1,20 = \textbf{525 k€}.

\textbf{Croissance.} Nominale : 630/500 − 1 = \textbf{26 \%}. Réelle : 525/500 − 1 = \textbf{5 \%}.

\textbf{Indice rebasé.} On divise par 108 puis on multiplie par 100 : \textbf{92,59 ; 100 ; 111,11}.

\textbf{Interprétation.} Le chiffre d'affaires en euros courants augmente de 26 \%, mais son pouvoir d'achat n'augmente que de 5 \% après correction de la hausse des prix. Le changement de base modifie l'affichage, pas la conclusion économique. Le déflateur doit rester pertinent pour les produits de l'entreprise.

\end{document}
